\documentclass[10pt,a4paper]{amsart}
\usepackage[margin=19mm]{geometry}
\usepackage{amsmath,amssymb,mathtools}
\usepackage[scr=rsfs,cal=euler]{mathalfa}
\usepackage{xfrac}
\usepackage[unicode,hidelinks,bookmarks=false]{hyperref}
\newcommand{\ie}{i.e.\ }
\newcommand{\cf}{cf.\ }
\newcommand{\resp}{respectively\ }
\newcommand{\Subsection}{\subsection}
\providecommand{\nobreakdash}{\nobreak\hbox{-}}
\setlength{\emergencystretch}{2em}
\multlinegap0pt
\usepackage{amssymb}
\usepackage{mathtools}
\usepackage{enumerate}
\usepackage{tikz}
\newtheorem{prop}{Proposition}
\newcommand{\N}{\mathbb{N}}
\newcommand{\Prob}{\operatorname{Prob}}
\newcommand{\bary}{\operatorname{Bar}}
\newcommand{\rC}{\operatorname{C}}
\newcommand{\rL}{\operatorname{ L}}
\newcommand{\supp}{\operatorname{supp}}
\title{Weyl groups and rigidity of \\ von Neumann algebras}
\author{Cyril Houdayer, Adrian Ioana}
\date{Source: arXiv:2508.08194v3 (CC BY 4.0)}
\begin{document}
\maketitle

\noindent\emph{Source and licence.} Weyl groups and rigidity of \\ von Neumann algebras. Version arXiv:2508.08194v3 (CC BY 4.0), distributed by arXiv under the Creative Commons Attribution 4.0 International licence, http://creativecommons.org/licenses/by/4.0/, as recorded in the arXiv licence metadata for this exact version and shown on the abstract page https://arxiv.org/abs/2508.08194v3. Full licence notice: \texttt{licenses/arxiv-2508.08194-CC-BY-4.0.txt}.\par\smallskip

\noindent\textit{Editorial scope.} Counted unit: the article's prop prop-ucp (source0.tex lines 407--412). The article's complete original proof of it is delivered with it (source0.tex lines 414--421, one \texttt{proof} environment of 8 lines). No mathematics is written, repaired or re-derived: the statement and the proof are the article's own lines, in the article's own notation, and no retained line is substituted or re-flowed.  No further source-local context is needed: every cross-reference of every retained line resolves inside the two delivered spans.  Nothing was omitted from any retained span, so the package carries no omission record.  No retained line contains a citation, so the wrapper carries no bibliography and no bibliography entry.  No retained line contains a figure environment, an image-inclusion command or drawing code, so no image is delivered and the package declares no asset dependency.  Editorial formatting only, in addition: an AMS \texttt{amsart} preamble that restates the article's own declarations and macros used by the retained lines, namely the environments \texttt{prop} on separate counters, together with the standard packages those lines need.  The retained environments print in this document as Proposition 1 in order of appearance, whereas the article prints the same items with the numbering of its own full document.  The complete native source of the article as distributed in the arXiv e-print for this exact version is bundled unchanged as \texttt{sources/183079/source0.tex}.
\begin{prop}\label{prop-ucp}
Let $L$ be a locally compact second countable group. Let $(X, \nu_X)$ be a standard probability $L$-space and $(Y, \nu_Y)$ a locally compact second countable Hausdorff topological $L$-space endowed with a fully supported Borel probability measure. 

To any $L$-equivariant normal ucp map $\Phi : \rL^\infty(Y, \nu_Y) \to \rL^\infty(X, \nu_X)$ corresponds an essentially unique $L$-equivariant measurable map $\beta : X \to \Prob(Y) : x \mapsto \beta_x$ such that $\bary(\beta_\ast \nu_X) = \nu_X \circ \Phi \prec \nu_Y$ and for $\nu_X$-almost every $x \in X$ and every $f \in \rC_0(Y)$, we have 
$$\beta_x(f) = \Phi(f)(x).$$
\end{prop}

\begin{proof} Let $\Phi : \rL^\infty(Y, \nu_Y) \to \rL^\infty(X, \nu_X)$ be a $L$-equivariant normal ucp map.

Firstly, we assume that $Y$ is compact. Since $\nu_Y$ is fully supported on $Y$, we may regard $\rC(Y) \subset \rL^\infty(Y, \nu_Y)$ as a unital $\rC^*$-subalgebra. Since $Y$ is a compact second countable Hausdorff topological space, $Y$ is metrizable and thus $\rC(Y)$ is $\|\cdot\|_\infty$-separable. Denote by $D \subset \rL^\infty(X, \nu_X)$ the $L$-invariant $\|\cdot\|_\infty$-separable unital $\rC^*$-subalgebra of $\rL^\infty(X, \nu_X)$ generated by the subspace $\Phi(\rC(Y))$. Observe that the action $L \curvearrowright D$ is $\|\cdot\|_\infty$-continuous. Denote by $Z$ the spectrum of $D$ and by $\nu_Z \in \Prob(Z)$ the Borel probability measure corresponding to the state $\nu_X|_{D} \in \mathfrak S(D)$. Then $Z$ is a compact metrizable topological $L$-space and we have  a $L$-equivariant measurable factor map $\pi : (X, \nu_X) \to (Z, \nu_Z)$. Then to the $L$-equivariant ucp map $\Psi : \rC(Y) \to \rC(Z) : f \mapsto \Phi(f)$ corresponds a unique $L$-equivariant continuous map $\alpha : Z \to \Prob(Y) : z \mapsto \alpha_z$ such that $\bary(\alpha_\ast \nu_Z) = \nu_Z \circ \Psi \prec \nu_Y$ and for every $z \in Z$ and every $f \in \rC(Y)$, we have $\alpha_z(f) = \Psi(f)(z)$. Then the $L$-equivariant measurable map $\beta : X \to \Prob(Y) : x \mapsto \alpha_{\pi(x)}$ does the job.

Secondly, we assume that $Y$ is noncompact. Since $\nu_Y$ is fully supported on $Y$, we may regard $\rC_0(Y) \subset \rL^\infty(Y, \nu_Y)$ as a $\rC^*$-subalgebra. Denote by $Z = Y \sqcup \{\infty\}$ the one-point compactification of $Y$, which is a compact metrizable space. Define the Borel probability measure $\nu_Z \in \Prob(Z)$ by $\nu_Z|_Y = \nu_Y$ and $\nu_Z(\{\infty\}) = 0$. Regard $\Prob(Y) \subset \Prob(Z)$ as a Borel subset. Then we have $\rL^\infty(Y, \nu_Y) = \rL^\infty(Z, \nu_Z)$. Let $\Phi : \rL^\infty(Y, \nu_Y) \to \rL^\infty(X, \nu_X)$ be an $L$-equivariant normal ucp map. Since the result holds for the compact metrizable space $Z$, there exists an essentially unique $L$-equivariant measurable map $\beta : X \to \Prob(Z) : x \mapsto \beta_x$ such that $\bary(\beta_\ast \nu_X) = \nu_X \circ \Phi \prec \nu_Z$ and for $\nu_X$-almost every $x \in X$ and every $f \in \rC(Z)$, we have $\beta_x(f) = \Phi(f)(x)$. It suffices to prove that $\beta_x(\{\infty\}) = 0$ for $\nu_X$-almost every $x \in X$.

We may find a nondecreasing sequence $f_n \in \rC_0(Y)$ of compactly supported functions such that $0 \leq f_n \leq 1$ for every $n \in \N$ and $f_n \to \mathbf 1_Y$ pointwise. For every $n \in \N$, denote by $K_n = \supp(f_n)$ the compact support in $Y$ of $f_n \in \rC_0(Y)$. Lebesgue's dominated convergence theorem implies that $f_n \to \mathbf 1_Y$ strongly in $\rL^\infty(Y, \nu_Y)$. Since $\Phi : \rL^\infty(Y, \nu_Y) \to \rL^\infty(X, \nu_X)$ is normal, we have $\Phi(f_n) \to \mathbf 1_X$ strongly in $\rL^\infty(X, \nu_X)$. Upon taking a subsequence, we may further assume that there is a conull Borel subset $X_0 \subset X$ for which $\beta_x(f_n) = \Phi(f_n)(x) \to 1$ for every $x \in X_0$. Then for every $x \in X_0$, we have $\beta_x(f_n) \leq \beta_x(K_n) \leq 1$ for every $n \in \N$ and so $\lim_n \beta_x(K_n) = 1$. Since $(K_n)_n$ is nondecreasing, letting $K = \bigcup_{n \in \N} K_n \subset Y$, for every $x \in X_0$, we have that $\beta_x(K) = 1$. Thus, for every $x \in X_0$, we obtain $\beta_x(\{\infty\}) \leq \beta_x(Z \setminus K) = \beta_x(Z) - \beta_x(K) = 0$.
\end{proof}

\end{document}
